\begin{table}[ht]
\centering
\caption{The capture ratio $r_{i,a}=S_{i,a}/G_{i,a}$ at anchor years, by geography tier.}
\label{tab:backcast_ratio}
\begin{adjustbox}{max width=\textwidth}
\begin{tabular}{lrrrrrr}
\toprule
& Anchors & p10 & p25 & Median & p75 & p90 \\
\midrule
Service-area tier & 814 & 0.842 & 1.068 & 1.530 & 3.169 & 11.530 \\
County tier & 306 & 0.0002 & 0.0005 & 0.0009 & 0.0029 & 0.0123 \\
\midrule
\multicolumn{7}{l}{\textit{Ratio drift: within-system $\max r / \min r$ across anchors (service-area tier)}} \\
Systems with $>1$ anchor & 235 & 1.00 & 1.04 & 1.15 & 1.44 & 2.19 \\
\bottomrule
\end{tabular}
\end{adjustbox}
\begin{minipage}{\textwidth}\footnotesize
\vspace{0.5em}
\textit{Notes:} $S$ is reported population served, $G$ the residential population inside the system's exposure polygon. For the service-area tier the ratio is interpretable and well behaved: it centers slightly above one, consistent with systems serving some demand beyond their resident population, and the median system's ratio moves only $1.15\times$ between its earliest and latest anchor, which is what the constant- and interpolated-$r$ assumptions require. For the county tier the denominator is an entire county's population, so $r\approx 0.001$ measures the utility's share of its county rather than a capture rate. Multiplying that ratio back by $G$ would simply return the reported figure while discarding the backbone, so county-tier systems instead take their reported population as the level and the county's growth as the trend.
\end{minipage}
\end{table}
